% !TeX root = ../main.tex
\section{Implementation and Software}
\label{sec:impl}

The game is written in Python~3 with the Pygame~2 library. The source is
open at:
\url{https://github.com/quyetops/Tactical-Formation-The-Board-Game}.

The source is organised into independent modules: \texttt{core/} (the rule
engine, board geometry, state), \texttt{ai/} (the Minimax and Paranoid
bots, policy), \texttt{ui/} (Pygame rendering, screens, widgets,
animation), and \texttt{tests/} (rule tests). The technical document
\texttt{docs/formal\_model.md} describes the formal mathematical model and
\texttt{docs/visual\_model.html} provides a visual companion.

\subsection{Mapping the formal model to source code}
\label{ssec:mapping}

The \texttt{core/} module directly implements the definitions in
Appendix~\ref{app:formal}, upholding the principle that \emph{the rule
engine depends on neither the interface nor the AI}:
\begin{itemize}[noitemsep,topsep=2pt]
  \item The state~\eqref{eq:state} is an immutable structure consisting of
        the board configuration $b$ (a cell $\to$ colour map) and the side
        to move.
  \item Move generation implements $\mathrm{Jumps}(x,b)$~\eqref{eq:moves}:
        for each piece of the side to move, it scans the direction set
        $D_{\kappa(x)}$ and filters by the three validity
        conditions~\eqref{eq:jump}.
  \item Sandwich detection implements the ray scan~\eqref{eq:clamp}: from
        the landing cell, along each capturing direction, it gathers the run
        of consecutive opponent pieces until it meets one of its own pieces
        (success), an empty cell, or the board edge (failure).
  \item Applying a move is the transition function
        $T$~\eqref{eq:t1}--\eqref{eq:t3}: jump, flip along the chosen
        direction, pass the turn---purely deterministic, creating no piece
        and deleting no piece (the invariant~\eqref{eq:conservation} is
        checked by an assertion that the total piece count $=28$ in
        \texttt{tests/}).
\end{itemize}
The bots in \texttt{ai/} call only those three operations (generate moves,
apply, evaluate): \texttt{minimax.py} implements
Algorithm~\ref{alg:minimax} for the 2-player mode, and
\texttt{paranoid.py} implements Algorithm~\ref{alg:paranoid} for the 3--4
player mode. Both are fully separated from the rule engine and can be
swapped for another algorithm without touching \texttt{core/}.

\begin{lstlisting}[language=bash, caption={Installing and running the game.}]
python -m venv .venv
source .venv/bin/activate   # Linux / macOS
pip install pygame
python main.py
\end{lstlisting}

The web version runs through \texttt{pygbag}~\citep{pygbag}:
\begin{lstlisting}[language=bash, caption={Running in the browser.}]
pip install pygbag
pygbag main.py
# Open the browser at http://localhost:8000
\end{lstlisting}
